% theme: gymnosperms_and_spore_plants
\MajorQuestion{植物のなかま（2）}
\SetRowSpaceQanda

\Instruction{マツの花と種子のでき方に関する次の問いに答えなさい。}
\QQ{マツのように、胚珠がむき出しになっている植物を\NamedBlank{gymnosperm}という。}{裸子植物}
\QQ{マツの新芽の先に集まり、りん片に胚珠がついている花を何というか。}{雌花}
\QQ{マツの新芽のつけ根に集まり、りん片で花粉をつくる花を\NamedBlank{malecone}という。}{雄花}
\QQ{\RefBlank{malecone}のりん片にあり、花粉をつくる袋を何というか。}{花粉のう}
\QQ{マツの花粉や種子は、おもに何によって運ばれるか。}{風}

\Instruction{種子をつくらない植物と植物全体の分類に関する次の問いに答えなさい。}
\QQ{種子をつくらず、小さな粒でふえる植物を\NamedBlank{sporeplant}という。}{胞子植物}
\QQ{\RefBlank{sporeplant}がなかまをふやすためにつくる小さな粒を\NamedBlank{spore}という。}{胞子}
\QQ{花を咲かせ、種子をつくってふえる植物を何というか。}{種子植物}
\QQ{胚珠が子房の中にある植物を何というか。}{被子植物}
\QQ{\RefBlank{gymnosperm}には、めしべの根元にある何が存在しないか。}{子房}

\Instruction{シダ植物のつくりに関する次の問いに答えなさい。}
\QQ{根・茎・葉の区別があり、葉の裏に袋をつけることが多い植物を\NamedBlank{fern}という。}{シダ植物}
\QQ{\RefBlank{fern}には、根・茎・葉の区別があるか、ないか。}{ある}
\QQ{地中を横にのびる、\RefBlank{fern}の茎を何というか。}{地下茎}
\QQ{\RefBlank{fern}の葉の裏にできる、\RefBlank{spore}をつくる袋を\NamedBlank{sporangium}という。}{胞子のう}
\QQ{\RefBlank{fern}の例を、イヌワラビ以外に1つ答えなさい。}{ゼンマイ（スギナ、ノキシノブ）}

\Instruction{コケ植物のつくりに関する次の問いに答えなさい。}
\QQ{根・茎・葉の区別がなく、体の表面全体から水を吸収する植物を\NamedBlank{moss}という。}{コケ植物}
\QQ{\RefBlank{moss}には、根・茎・葉の区別があるか、ないか。}{ない}
\QQ{\RefBlank{moss}は体のどの部分から水を吸収するか。}{体の表面全体}
\QQ{\RefBlank{moss}の体を地面に固定する、根のような部分を何というか。}{仮根}
\QQ{\RefBlank{moss}で、\RefBlank{sporangium}ができる株を何というか。}{雌株}
